% !TEX root = ../oddp.tex

\subsection{The bar resolution}\label{ss:milnor} Recall the that the Milgram construction applied to a group $G$, yields a contractible augmented simplicial set $\EG$ with a free $G$-action:
\begin{align*}
	(\EG)_k &= G^{k+1}\\
	\d_i(g_0\dots g_k) &= (g_0\dots\hat{g}_i\dots g_k) \\
	\s_i(g_0\dots g_k) &= (g_0\dots g_i,g_i\dots g_k)
\end{align*}
We can regard $\EC_r$ as the union $\bigcup_n \chains(\Cyc_r)^{\ot n}$ (this filters $\EC_r$ by degree). Then, the inclusion of $\Cyc_r$ as the vertices of $\asimplex^{r-1}$ induces a map $\chains(\EC_r) \to \bigcup_n\chains(\asimplex^{r-1})^{\ot n}$ that sends a tuple $(a_0,a_1,\dots,a_k)$ to the generator $[a_0]\ot \ldots \ot [a_k]$. The composition
\[
\chains(\EC_r) \lra \bigcup_n\chains(\asimplex^{r-1})^{\ot n} \lra \Waug{r}
\]
is an equivariant chain map. We can describe the dual of this map as follows: A generator $(-1)^{s}\cdot c\cdot (a_0,\ldots,a_{q-1})$ is \emph{admissible} if 
\begin{itemize}
	\item The generator $(-1)^{s'}\cdot (c\cdot \tilde{r}!) [a_0,\ldots,a_{q-r}]$ is admissible and $[a_{q-r},\ldots,a_{k-1}]$ does not contain repetitions and $(-1)^{s'+s}$ computes the parity of this permutation.
	\item $q<r$, then $s$ is the sign of the permutation that orders $[a_0,\ldots,a_{q-1}]$, $c$ equals $\frac{\varphi(q)}{\tilde{r}!}$, the number $r-1$ does not belong to the word, and the word obtained after ordering it starts with an even entry and alternates even and odd entries.
\end{itemize}

%
 %This composition has a description analogous to the one used to define $\Psiom$, using the following suspension map:
%\[
%S(a_0,\dots,a_k) = \begin{cases}
	%(-1)^{\sign{\perm}} (a_{0},\dots,a_{k-r+1}) & \text{ if $k \geq r-1$} \\
	%\emptyset & \text{if $k = r-2$ and $a_i\neq a_j$ for all $i,j$} \\
	%0 & \text{otherwise.}
%\end{cases}
%\]
%where $\pi$ is the permutation that orders $(a_{k-r+1},\dots,a_{k})$. The image of $e_{-q}^{\vee}$ is the sum of all words $A = (a_0,\dots,a_{q-1})\in \cochains[-q] \EC_r$ such that
%\begin{itemize}
%\item none of the subwords of the form $(a_{q-1-\ell(r-1)},a_{q-1-(\ell-1)(r-1))}$ of length $r$ contains repeated elements,
%\item the only subword of the form $(a_0,\dots,a_{q-1-\ell(r-1)})$ of length smaller than $r$ becomes, after having been ordered, an increasing sequence that alternates even and odd entries different from $r-1$, starting with an even entry.
%\end{itemize}
%each word has the sign needed to order each of the subwords mentioned above, and the coefficient $\frac{\varphi(m)!}{(\tilde{r}!)^j}$ where $m$ is the length of the last of the subwords mentioned, and $j$ is the total number of subwords and $\varphi(m) = \lfloor\frac{r-m-1}{2}\rfloor$.
\begin{example} If $r=3$,
\begin{align*}
	\Psi^\vee(e_{-1}^\vee) &= (0)
	\\
	\Psi^\vee(e_{-2}^{\vee}) &= (0,1)-(1,0)
	\\
	\Psi^\vee(e_{-3}^{\vee}) &= (0,1,2) - (0,2,1)
\\
	\Psi^\vee(e_{-4}^{\vee}) &= (0,1,2,0) - (0,1,0,2) + (1,0,2,1) - (1,0,1,2)
\end{align*}
if $r= 5$,
\begin{align*}
	\Psi^\vee(e^\vee_{-1}) &= \frac{1}{2}\left((0) + (2)\right)
	\\
	\Psi^\vee(e^\vee_{-2}) &= \frac{1}{2}\left((0,1) - (1,0) + (0,3) - (3,0) + (2,3) + (3,2)\right)
	\\
	\Psi^\vee(e^\vee_{-3}) &= \frac{1}{2}\left((0,1,2) -(0,2,1) + (1,2,0) - (1,0,2) + (2,0,1) - (2,1,0)\right)
\end{align*}
$\Psi^\vee(e^\vee_{-4})$ is the signed sum of all permutations of $(0,1,2,3)$, so it has $24$ summands, again with coefficient $\frac{1}{2}$.
\end{example}

This map is different from the one obtained in \cite[Prop.~6.16]{brumfiel2023explicit}, and in fact that map cannot be used to construct $\Psiom$ following the recipy of Definition \ref{def:psiom}.

\subsection{The join and the reduced join} Observe that the map $\psi$ constructed factors through the Milgram resolution:
\[
	\chains(\partial\asimplex^{r-1})^{\ot \bullet+1}
	\lra
	\bigcup_n \chains(\partial \simplex^{r-1})^{\ot n}_+
	\lra
	\Per{r}
\]
where the middle chain complex has been augmented. For example, the image of $\tau_0\ot \emptyset \ot \tau_1$ is the same as the image of $\tau_0\ot\tau_1\ot\emptyset$.


\subsection{The even prime}\label{ss:even} We have only constructed a map $\phi$ when $r$ is odd, but if $r= 2$ and $R=\bF_2$, then $\Per{r}$ and $\Waug{r}$ are isomorphic and we can define an isomorphism $\phi\colon \Per{r}\to \Waug{r}$ sending the generator $\sus{q-1}[0]$ to $e_q$.

Regarding the map $\psi$, observe that $\partial \simplex^{r-1}$ is isomorphic to $\Cyc_r$, thus the union $\bigcup_n \chains(\partial \simplex^1)^{\ot n}$ is isomorphic to the unnormalized chain complex of the Milgram resolution $\EC_2$, while $\Per{r}$ is isomorphic to the normalized chain complex of $\EC_2$. The map $\psi\colon \chains(\partial \asimplex^1)^{\ot (n+1)} \lra \Per{r}$ is the composition of
\[
	\chains(\partial \asimplex^1)^{\ot (n+1)}
	\to
	\uchains(\EC_r)
	\to
	\chains(\EC_r)
	\overset{\cong}{\to}
	\Per{r}
\] 
where the first map forgets the empty factors, the second map is the projection to the normalised chains and the last map sends a normalised chain $[u_0,u_1,\ldots,u_{q-1}]$ to $\sus{q-1}[u_0]$. 

Under this identification, the coface map $\d^i$ of $\chains(\partial\simplex^{r-1})^{n+1}$ (resp.\ $\chains(\partial\asimplex^{r-1})^{n+1}$) sends a sequence $[a_0\ot \ldots\ot a_n]$ of $[0]$'s and $[1]$'s (resp.\ and $\emptyset$) to the sum of the two sequences that result from adding a $[0]$ or a $[1]$ in position $i$ and applying $\rho^{-1}$ to the terms in position $\geq i$. The coface map $\d^i$ of $\Per{r}$ sends $\sus{q}[0]$ to $\sus{q+1}[0]$ and $\sus{q}[1]$ to $\sus{q+1}[1]$. Thus $\psi$ commutes with $\d^i$, and therefore is a cosemi-simplicial chain map.
\begin{remark}
	The combinatorics in pages 13 and 14 \federico{check pages, I don't have access to the published version} of \cite{medina2021fast_sq} translate here to the fact that the map from the unnormalized chains to the normalized chains is well-defined, i.e., to the fact that the degenerate simplices generate a subcomplex of the unnormalized chain complex. To add more detail: Medina-Mardones needs to check that the map $\alpha^\vee\circ \beta^{\vee}\circ \psi^\vee\circ \phi^\vee$ (denoted $\nabla^{\binom{n}{q}}$) is a cosemi-simplicial chain map, though most of the trouble is just proving that it is a chain map.
	
	As we have seen in the previous subsection, the map $\psi$ factors as $\mathrm{red}\circ \psi_{\mathrm{red}}$ and we have
	\[
		e_q^{\vee}\overset{\phi^\vee}{\longmapsto}
		\susp{q}(0)\overset{\psi^\vee_{\mathrm{red}}}{\longmapsto}
		(0)\ot(1)\ot(0)\overset{q}{\ldots} 
	\]
	Now, the map $\mathrm{red}$ takes this last generator and inserts $n+1-q$ instances of $\emptyset$ in all possible ways, for example, if $q=2$ and $n=3$, then we obtain the following summands:
	\begin{align*}
		&(0)\ot(1)\ot\emptyset\ot\emptyset
		&
		&(0)\ot\emptyset\ot(1)\ot\emptyset
		&
		&(0)\ot\emptyset\ot\emptyset\ot(1)
		\\
		&\emptyset\ot(0)\ot(1)\ot\emptyset
		&
		&\emptyset\ot(0)\ot\emptyset\ot(1)
		&
		&\emptyset\ot\emptyset\ot(0)\ot(1)
	\end{align*}
	Each of these summands can be indexed by the positions of the non-empty entries
	\begin{align*}
		&(0,1)&& (0,2)&&(0,3)
		\\
		&(1,2)&& (1,3) &&(2,3)
	\end{align*}
	These are the elements $U\in \cP_q(n)$ in Medina-Mardones' notation. The map $\beta^\vee$ sends these summands to
	\begin{align*}
		&(0)\ot(0)\ot\emptyset\ot\emptyset
		&
		&(0)\ot\emptyset\ot(1)\ot\emptyset
		&
		&(0)\ot\emptyset\ot\emptyset\ot(0)
		\\
		&\emptyset\ot(1)\ot(1)\ot\emptyset
		&
		&\emptyset\ot(1)\ot\emptyset\ot(0)
		&
		&\emptyset\ot\emptyset\ot(0)\ot(0)
	\end{align*}
	and the map $\alpha$ sends these summands to
	\begin{align*}
		&(0,1)\ot \emptyset
		&
		&(0)\ot (2)
		&
		&(0,3)\ot \emptyset
		\\
		&\emptyset\ot(1,2)
		&
		&(3)\ot (1)
		&
		&(2,3)\ot\emptyset
	\end{align*}
	which correspond to the decomposition $U^0\ot U^1$ of each $U$ in Medina-Mardones' notation. 

	The cancellations taking place in these pages are due to the fact that $\psi_{\mathrm{red}}^\vee \circ \phi^\vee(e_q^\vee)$ in $\cochains(\partial\simplex^{r-1})^{\ot n+1}\cong \ucochains(\EC_2)$ is a normalised cochain, thus its coboundary is again a normalised cochain. For example, in $\cochains(\asimplex^{n})^{\ot r}_{\interior}$ the coboundary of $(0)\ot (2)$ has a summand $(0,1)\ot (2)$ that also appears in the coboundary of $(0,1)\ot\emptyset$, thus they cancel out. If we translate this back to the chain complex of the reduced join $\cochains(\partial\simplex^{r-1})^{\ot 4}$, we see that $(0)\ot (2)$ and $(0,1)\ot\emptyset$ are summands in the image of $(0)\ot (1)$ and that $(0,1)\ot (2)$ is a summand in the image of $(0)\ot (1)\ot(1)$. Now, in the coboundary of $(0)\ot (1)$ the summand $(0)\ot (1)\ot(1)$ appears twice: one comes from inserting a $(1)$ in the middle position and another from inserting a $(1)$ in the right position, thus they cancel out.
\end{remark}